\scp{Possible incursion and migration scenarios}{15-06}
In this section we investigate some possibilities regarding the
prehistory of eastern Indonesian
and Timor-Leste based on the revised classification proposed in this work.
Before we proceed, we want to emphasise the preliminary nature of these proposals.
They are intended to spark ideas about reasonable possibilities,
and to underscore problems with assuming a single northeast-to-southwest
movement proposed by \citet[98]{Blust2008} for the breakup of PCEMP.

Phonologically conservative features (retentions) in many subgroups
imply that the Austronesian languages spoken when they entered
and spread through the target region diverged very little from PMP.
Conservative features that have been discussed in previous chapters
include: retention of the nasals; retention of the labials,
including retaining distinctions between *-mp-,
*-mb-, *ma-p and *ma-b; retention of vowel-glide sequences; retention
of PMP *ə; retention of PMP *q
and *y, (and perhaps *h); retention of many word-final consonants;
retention of medial *NC clusters;
and retention in the region of both parts of other medial *CC clusters.
Retention of the four-vowel system of PMP explains the data better
than does asserting the development in a putative proto-language
of mid-vowels *e and *o.
(It is important to keep in mind that retentions do not define
subgroups; exclusively shared innovations are what define subgroups.)

The scattered distribution of many conservative features implies
the the Austronesian-speaking peoples either spread very quickly
through the region, or conservative varieties of Austronesian languages entered the
region in and from several different places or at different times.
Or all of the above.
The diversity, distribution and discontinuities of many features
found in the region imply there was not just one variety of a
single Austronesian language, but that a number of languages more diverse than a dialect chain
or linkage spread through the region,
probably from different origin points
and in different directions.

The starting baseline prior to the arrival of Austronesian-speaking
peoples is the presence of pre-existing populations scattered
throughout the region at least as far west as Sumbawa already
speaking various Papuan languages,
most of which were probably typologically SOV with preposed possessors,
an alienable-inalienable distinction in their possessive constructions,
post-positions, and post-verbal or clause final negation.

We note in near Oceania where significantly more archaeology
has been done in relation to the expansion of the Lapita cultural
complex than is comparable anywhere in eastern Indonesia
and Timor-Leste, it has been found that there are places that
are “best explained as a leapfrog rather than a wave-of-advance
movement out of New Guinea into Remote Oceania” \citep{Sheppard2019}.
In other words people in different regions got off different
boats without stopping along the way where others stopped.
This, along with contact, provides a likely model for many of
the complexities and discontinuities we find in the Austronesian languages of
eastern Indonesia and Timor-Leste.

It appears that speakers of several different Austronesian languages
leap-frogged and settled in different places,
each having unique patterns of contact with the Papuan languages
in their region, and perhaps with other groups of Austronesian speakers.
Several different subgroups developed independently through shared
innovations, complicated by unique and different patterns of intense contact
and bilingualism that continued over long periods of time -- probably millennia.
At the same time there were a number of useful words that diffused
unevenly throughout the region,
often blind to whether a language was Austronesian or Papuan.

The diversity of subgroups in the region do not support a single
movement of Austronesian languages into the region from a single
point or in a single direction.
For several subgroups (e.g.
\tsc{Timor-Babar}) the current evidence does not allow us to
identify an origin point or direction of movement,
while for others (e.g.
\tsc{Bima-Lembata}) there is better evidence for a direction of movement.
We also note that further refinement/revision to the subgroups
we have proposed in this book may provide additional information
on direction of movement.
Below we examine evidence for the following possible incursion
and migration scenarios into
and within the target region:

\begin{itemize}
		\item	\it{Austronesian incursion} (\it{major}): \tsc{Bima-Lembata},
					from west-to-east across the southern islands.
		\item	\it{Austronesian incursion} (\it{minor}): \tsc{Taliabu},
					from west-to-east from central Sulawesi into central Maluku.
		\item	\it{Austronesian incursion} (\it{possible}): \tsc{Sula-Buru}?,
					\tsc{Ambon-Seram}?, \tsc{Seram-Tanimbar-Bomberai}?, from northeast to southwest.
					From Halmahera-Bird’s Head through central Maluku,
					southern Maluku, to the Lesser Sundas.
		\item	\it{Papuan incursion} (\it{secondary}): from east-to-west.
					Probably prior to the arrival of the Austronesians.
		\item	\it{Papuan migrations} (\it{secondary}): \tsc{Alor-Pantar},
					Oirata, from the Timor mainland to Alor-Pantar and to Kisar.
		\item	\it{Austronesian incursion}/\it{migration}: \tsc{Timor-Babar}.
		\item	\it{Austronesian incursion}/\it{migration}: \tsc{Central Timor}.
		\item	\it{Austronesian incursion}/\it{migration}: \tsc{Aru}.
		\item	\it{Austronesian incursion}/\it{migration}: \tsc{Seram-Tanimbar-Bomberai}.
		\item	\it{Austronesian incursion}/\it{migration}: \tsc{South Halmahera-West New Guinea}, east-to-west.
					\citet[142]{Kamholz2014} describes the spread of the \tsc{Raja
					Ampat-South Halmahera} node as moving from southern Cenderawasih
					Bay around the coastal Bird’s Head
					and Raja Ampat Islands to South Halmahera.
					This development within SHWNG is part of the wider picture in
					the region, but beyond the scope of this study.
\end{itemize}

\frf{15-10} shows some of these possible migration scenarios
to provide a visual reference for the discussion that follows
in subsequent sections.

\begin{figure}[h]\centering
	\includegraphics[width=\columnwidth]{migrations.pdf}
	\fcp{Possible incursion and migration scenarios}{15-10}
\end{figure}